\documentclass[11pt,oneside]{article}
\usepackage[a4paper,margin=27mm,headheight=15pt]{geometry}
\usepackage[T1]{fontenc}
\usepackage[utf8]{inputenc}
\usepackage{libertinus-type1}
\usepackage{microtype}
\usepackage{amsmath,amssymb,amsthm,mathtools,mathrsfs}
\usepackage{bm}
\usepackage{booktabs,longtable,array,tabularx}
\usepackage{graphicx,xcolor,enumitem,fancyhdr,titlesec,listings,xurl}
\usepackage{hyperref}
\usepackage[nameinlink,noabbrev]{cleveref}
\definecolor{linkblue}{RGB}{25,75,135}
\definecolor{shadegray}{RGB}{246,247,249}
\definecolor{rulegray}{RGB}{195,201,208}
\hypersetup{colorlinks=true,linkcolor=linkblue,citecolor=linkblue,urlcolor=linkblue,
 pdftitle={Second Order Critical Complements in Fabius Conditioning},
 pdfauthor={Research note prepared with OpenAI for Vladimir Reshetnikov},
 pdfsubject={Uniform overlap corrections and a strict comparison of observation masks}}
\pagestyle{fancy}\fancyhf{}
\fancyhead[L]{\small Fabius critical complements}
\fancyhead[R]{\small Research note}
\fancyfoot[C]{\thepage}
\renewcommand{\headrulewidth}{0.4pt}
\titleformat{\section}{\Large\bfseries}{\thesection}{0.7em}{}
\titleformat{\subsection}{\large\bfseries}{\thesubsection}{0.7em}{}
\setlength{\emergencystretch}{2em}
\setlist[itemize]{leftmargin=2em,itemsep=0.25em,topsep=0.4em}
\setlist[enumerate]{leftmargin=2.3em,itemsep=0.25em,topsep=0.4em}
\newtheorem{theorem}{Theorem}[section]
\newtheorem{proposition}[theorem]{Proposition}
\newtheorem{lemma}[theorem]{Lemma}
\newtheorem{corollary}[theorem]{Corollary}
\theoremstyle{remark}\newtheorem{remark}[theorem]{Remark}
\newcommand{\E}{\mathbb E}
\newcommand{\Pp}{\mathbb P}
\newcommand{\R}{\mathbb R}
\newcommand{\N}{\mathbb N}
\newcommand{\dd}{\,\mathrm d}
\newcommand{\Law}{\mathcal L}
\newcommand{\eps}{\varepsilon}
\newcommand{\TV}{d_{\mathrm{TV}}}
\newcommand{\ind}{\mathbf 1}
\newcommand{\OO}{\mathcal O}
\newcommand{\HH}{\mathcal H}
\newcommand{\bracket}[1]{\left[#1\right]_{y_-}^{y_+}}
\numberwithin{equation}{section}
\title{\bfseries Second Order Critical Complements\\in Fabius Conditioning\\[0.4em]
\large Uniform overlap corrections and a strict comparison of two observation masks}
\author{Research note prepared for Vladimir Reshetnikov\\\small Developed with OpenAI}
\date{1 October 2026}

\begin{document}
\maketitle
\begin{abstract}
Two recent ProveIt manuscripts study near-total observation of a geometrically
weighted uniform series under a small-value constraint. They omit different
bulk coordinates. We derive a uniform expansion of their overlap when the
number of omitted bulk coordinates is comparable with the logarithm of the
bulk size. The expansion includes the phase-dependent constant and the next
inverse-complement coefficient, with a quantitative remainder. It answers the
second-order crossover question in the prefix-observation manuscript and the
logarithmic slice of the other manuscript's growing-complement question.
The two masks have the same leading crossover but different next constants:
hiding the earliest bulk coordinates gives strictly greater overlap than
hiding the latest ones, eventually. An elementary signed Gamma convolution
argument proves the statements without importing a uniform conditioning
theorem. These are manuscript-level proofs, not a proof-assistant
formalization or a certification of worldwide novelty.
\end{abstract}

\section{Scope and relation to the source questions}
This note examines the repository snapshot
\nolinkurl{63a7a325109ba611a1816b61dfd0eb072b896a7a}.
The article in \emph{Critical Complements Sharp Information Loss}
\cite{repo-late}, Section 10, asks for the complete next-order overlap
expansion when $m\sim c\log n$, including the geometric phase and integer
rounding. Its Theorem 2.2 already proves the leading equivalent for every
$m=o(n)$. Here \cref{thm:main,cor:loglog} give a finer expansion in the
critical window, including an additional $1/m$ coefficient.

The article in \emph{Critical Complements Fabius Conditioning}
\cite{repo-early}, Section 13, instead observes a block obtained by deleting
the earliest $m$ bulk coordinates. It leaves growing $m$ open. The block
is not the prefix in \cite{repo-late}. We prove the same expansion for
this second mask, with a different analytic factor, and establish the
strict comparison in \cref{cor:comparison}. We do not claim that the
leading theorem for one mask previously solved the question for the other.

The probability representation of the Fabius function is classical
\cite{arias}. Exponential tilting, Gamma asymptotics, and precise-deviation
methods are also established tools; compare \cite{fmnn,dlmf}.
The proposed contribution here is the explicit clipped-density expansion
and its application to these two specified masks. The literature check was
targeted, not exhaustive. No famous conjecture, publication priority, or
Lean verification is claimed.

\section{The geometric model and the two masks}
Fix $q\in(0,1)$ and $\rho>0$. Let
\begin{equation}
 S_q=\sum_{j\ge1}q^jU_j,\qquad U_j\overset{\mathrm{iid}}\sim
 \operatorname{Unif}[0,1],\qquad t_n=\rho q^{-(n+1)}.
\end{equation}
Write $P$ for the product law, $T_n=t_nS_q$, and define
\begin{equation}
 \frac{\dd Q_n}{\dd P}=\frac{e^{-T_n}}{\E_P e^{-T_n}},\qquad
 \mu_n=\E_{Q_n}T_n,\qquad P_n=P(\,\cdot\mid T_n\le\mu_n).
 \label{eq:model}
\end{equation}
The constraint is an inequality under the original law. Under $Q_n$,
the scaled coordinates are independent with truncated exponential density
\begin{equation}
 p_a(x)=\frac{e^{-x}}{1-e^{-a}}\ind_{(0,a)}(x),\qquad
 a=\rho q^{j-(n+1)}.
\end{equation}
The first $n$ coordinates form the bulk. The boundary tail starts at
coordinate $n+1$, of cap $\rho$. Set
\begin{equation}
 a_r=\rho q^{-r}\quad(r\ge1),\qquad
 R_0=\sum_{\ell\ge0}Y_\ell,\qquad Y_\ell\sim p_{\rho q^\ell}
 \quad\hbox{independently}.
 \label{eq:tail}
\end{equation}
Thus $0\le R_0\le\rho/(1-q)$. For $1\le m<n$, define two observed sets:
\begin{align}
 B^{\mathrm{late}}_{n,m}&=\{1,\ldots,n-m\},
 &B^{\mathrm{early}}_{n,m}&=\{m+1,\ldots,n\}.
 \label{eq:masks}
\end{align}
The superscript names the hidden bulk coordinates. Both masks hide the
whole infinite boundary tail. If $P^j_{n,m},Q^j_{n,m}$ are the corresponding
vector marginal laws, let
\begin{equation}
 \OO^j_{n,m}=1-\TV(P^j_{n,m},Q^j_{n,m}),\qquad
 \eps_n=(2\pi n)^{-1/2},\qquad j\in\{\mathrm{early},\mathrm{late}\}.
\end{equation}
Total variation is $\sup_A|P(A)-Q(A)|$, so overlap is
$\int\min(\dd P,\dd Q)$.

The phase factors needed below are
\begin{align}
 H_0(s)&=\E e^{sR_0},\label{eq:H0}\\
 \HH_{\mathrm{early}}(s)&=H_0(s),\label{eq:HE}\\
 \HH_{\mathrm{late}}(s)&=H_0(s)
 \prod_{r\ge1}\frac{1-e^{-(1-s)a_r}}{1-e^{-a_r}},\qquad s<1.
 \label{eq:HL}
\end{align}
The product converges locally uniformly, with every derivative, in the
complex half-plane $\Re s<1$. Both factors are strictly positive for
real $s<1$ and equal one at zero. If desired, $H_0$ itself has the
explicit product
\begin{equation}
 H_0(s)=\prod_{\ell\ge0}
 \frac{1-e^{-(1-s)\rho q^\ell}}
 {(1-s)(1-e^{-\rho q^\ell})},
 \label{eq:H0product}
\end{equation}
with each removable value interpreted continuously.

The first source \cite{repo-early} uses $t_n=\rho_{\rm first}q^{-n}$.
The exact phase match to our convention is $\rho_{\rm first}=\rho/q$.
No asymptotic reindexing is hidden in this comparison.

\section{The uniform expansion}
Let
\begin{equation}
 I(y)=y-1-\log y,\quad s(y)=I'(y)=1-\frac1y,\quad
 d=\frac12\log\frac nm,\quad z=\frac dm.
 \label{eq:rootsdata}
\end{equation}
The two roots of $I(y)=z$ are denoted $y_-<1<y_+$.
For a positive analytic factor $\HH$ from \eqref{eq:HE} or \eqref{eq:HL},
define
\begin{align}
 L(y)&=\log\HH(s(y)),\label{eq:L}\\
 A(y)&=-\frac1{12}-\frac{\HH''(s(y))}{2y^2\HH(s(y))},\label{eq:A}\\
 F(y)&=\frac{L(y)+1}{s(y)},\label{eq:F}\\
 B(y)&=\frac{A(y)}{s(y)}+
       \frac{(L(y)+1)L'(y)}{s(y)^2}
       -\frac{L(y)^2/2+L(y)+1}{y^2s(y)^3}.
 \label{eq:B}
\end{align}
Here $\HH''$ is differentiated in $s$ and $L'$ in $y$;
in particular $L'(y)=\HH'(s(y))/(y^2\HH(s(y)))$.
We write $[G]_{y_-}^{y_+}=G(y_+)-G(y_-)$.

\begin{theorem}[Critical crossover through the inverse-complement term]
\label{thm:main}
Fix $q,\rho$ and $0<c_0<c_1<\infty$. Uniformly over integers $m$ such that
$c_0\log n\le m\le c_1\log n$, as $n\to\infty$, one has
\begin{equation}
 \frac{\OO^j_{n,m}}{\eps_n}
 =m(y_+-y_-)+[F_j]_{y_-}^{y_+}
       +\frac1m[B_j]_{y_-}^{y_+}
       +O\!\left(\frac1{m^2}+\frac{m^3}{n}\right).
 \label{eq:main}
\end{equation}
The factors in \eqref{eq:HE}--\eqref{eq:HL} determine $F_j,B_j$.
All roots use the actual integer $m$ and $d=\frac12\log(n/m)$.
Consequently the same formula with the $1/m$ term omitted has
remainder $O(1/m)$.
\end{theorem}

The geometric phase first appears in the bounded correction $[F_j]$;
the leading term forgets it. The two exterior tails of the clipped density
each contribute a nonzero part of this correction. Keeping only the
distance between likelihood crossings would omit $[1/s]_{y_-}^{y_+}$.

\subsection{The logarithmic and rounding terms}
For $c>0$, let $Y_\pm(c)$ solve $I(Y)=1/(2c)$ and set
\begin{align}
 W(c)&=Y_+(c)-Y_-(c),&
 C(c)&=cW(c),\\
 D(c)&=\frac1{s(Y_+(c))}-\frac1{s(Y_-(c))},&
 K_j(c)&=[F_j]_{Y_-(c)}^{Y_+(c)}.
\end{align}
\begin{corollary}[Explicit logarithmic correction]
\label[corollary]{cor:loglog}
Put $\ell=\log n$ and $c_n=m/\ell$. Uniformly for $c_n\in[c_0,c_1]$,
\begin{equation}
 \frac{\OO^j_{n,m}}{\eps_n}
 =\ell C(c_n)-\frac12D(c_n)\log\ell
       +K_j(c_n)-\frac12D(c_n)\log c_n
       +O\!\left(\frac{(\log\ell)^2}{\ell}
                         +\frac{\ell^3}{n}\right).
 \label{eq:loglog}
\end{equation}
If $m=c\ell+\beta_n$ for fixed $c>0$ and bounded $\beta_n$, then
\begin{align}
 \frac{\OO^j_{n,m}}{\eps_n}
 ={}&\ell C(c)-\frac12D(c)\log\ell
       +\beta_n\left(W(c)-\frac{D(c)}{2c}\right)\nonumber\\
 &+K_j(c)-\frac12D(c)\log c+o(1).
 \label{eq:rounding}
\end{align}
Thus rounding $m=\lfloor c\log n\rfloor$ generally produces an
oscillating bounded correction. The assumption $m\sim c\log n$ alone
does not permit replacing $c_n$ by $c$ in \eqref{eq:loglog} to constant
accuracy.
\end{corollary}
\begin{proof}
If $\Delta(z)=y_+(z)-y_-(z)$, implicit differentiation gives
$\Delta'(z)=1/s(y_+)-1/s(y_-)$. Since
$z=1/(2c_n)-\log(c_n\ell)/(2c_n\ell)$, Taylor's theorem applied
on a compact positive interval gives
\[
 m\Delta(z)=\ell C(c_n)-\tfrac12D(c_n)\log(c_n\ell)
                  +O((\log\ell)^2/\ell).
\]
Replacing the arguments of $F_j$ costs $O(\log\ell/\ell)$.
This proves \eqref{eq:loglog}. Finally
$C'(c)=W(c)-D(c)/(2c)$, which gives \eqref{eq:rounding}.
\end{proof}

\subsection{A strict distinction between the two masks}
\begin{corollary}[Hiding earlier coordinates gives greater overlap]
\label[corollary]{cor:comparison}
Let $m/\log n\to c\in(0,\infty)$. Put
\begin{equation}
 \Psi(s)=\sum_{r\ge1}
  \log\frac{1-e^{-(1-s)a_r}}{1-e^{-a_r}}.
 \label{eq:Psi}
\end{equation}
Then, with $s_\pm=s(Y_\pm(c))$,
\begin{equation}
 \frac{\OO^{\mathrm{late}}_{n,m}-\OO^{\mathrm{early}}_{n,m}}{\eps_n}
 \longrightarrow \frac{\Psi(s_+)}{s_+}-\frac{\Psi(s_-)}{s_-}<0.
 \label{eq:ordering}
\end{equation}
In particular $\OO^{\mathrm{late}}_{n,m}<\OO^{\mathrm{early}}_{n,m}$
for all sufficiently large $n$ along that sequence.
\end{corollary}
\begin{proof}
The leading root term and the $1/s$ tail terms cancel in the difference.
Since $\log\HH_{\mathrm{late}}-\log\HH_{\mathrm{early}}=\Psi$,
\cref{thm:main} gives the stated limit. Termwise differentiation is valid,
and
\[
 \Psi''(s)=-\sum_{r\ge1}
 \frac{a_r^2e^{-(1-s)a_r}}{(1-e^{-(1-s)a_r})^2}<0,
 \qquad \Psi(0)=0.
\]
Strict concavity implies that the secant slope on $[s_-,0]$ is strictly
larger than that on $[0,s_+]$. These slopes are
$\Psi(s_-)/s_-$ and $\Psi(s_+)/s_+$, respectively.
\end{proof}

\section{A clipped signed Gamma transfer theorem}
The following analytic statement isolates the calculation. It is deliberately
stated with hypotheses that can be checked directly for both masks.
Write
\[
 \gamma_k(x)=\frac{e^{-x}x^{k-1}}{\Gamma(k)}\ind_{(0,\infty)}(x),
 \qquad \|\eta\|_\theta=\int_0^\infty e^{\theta u}|\eta|(\dd u).
\]
\begin{lemma}[Density and endpoint tail expansions]
\label[lemma]{lem:saddle}
Let $\eta_m$ be real finite signed measures supported in $[0,\infty)$,
of mass one. Suppose, for every $\theta<1$,
\[
 \sup_m\|\eta_m\|_\theta<\infty,
 \qquad \|\eta_m-\eta\|_\theta=O(m^{-2}),
\]
and put $\HH(s)=\int e^{su}\eta(\dd u)$. Suppose $\HH(s)>0$ for real
$s<1$ and $g_m=\gamma_{m+1}*\eta_m$ is a nonnegative probability density.
Uniformly for $y$ in compact subsets of $(0,\infty)$,
\begin{equation}
 g_m(my)=\frac{e^{-mI(y)}\HH(s(y))}{\sqrt{2\pi m}}
        \left(1+\frac{A(y)}m+O(m^{-2})\right).
 \label{eq:saddle}
\end{equation}
Here $A$ is \eqref{eq:A}. If the convergence of $\eta_m$ above is
faster than every inverse power, the hypothesis is in particular satisfied.
\end{lemma}
\begin{proof}
The exact Gamma ratio is
\begin{equation}
 \frac{\gamma_{m+1}(my-u)}{\gamma_{m+1}(my)}
 =e^u\left(1-\frac{u}{my}\right)^m\ind_{\{u<my\}}.
 \label{eq:gammaratio}
\end{equation}
For $y$ in a fixed compact interval and $0\le u\le my/2$, expanding
$\log(1-u/(my))$ and then the exponential gives
\begin{equation}
 e^{s(y)u}\left(1-\frac{u^2}{2my^2}\right)
 +O\!\left(\frac{(u^3+u^4)e^{s(y)u}}{m^2}\right).
 \label{eq:ratioexp}
\end{equation}
For completeness, the extra negative exponent beyond its quadratic term is
bounded by $Cu^3/m^2$ on this interval. The inequalities
$|e^{-v}-1+v|\le v^2/2$ and $1-e^{-w}\le w$ for $v,w\ge0$
prove \eqref{eq:ratioexp} with an absolute bound of the displayed form.
On $u>my/2$, both the exact ratio and the approximating terms are
exponentially negligible after integration: choose
$\theta\in(\max_y s(y),1)$ and use the weighted variation bound.
Thus integration against the signed measure gives
\[
 \int\frac{\gamma_{m+1}(my-u)}{\gamma_{m+1}(my)}\eta_m(\dd u)
 =\HH(s(y))-\frac{\HH''(s(y))}{2my^2}+O(m^{-2}).
\]
The $O(m^{-2})$ convergence hypothesis permits replacement of all needed
weighted moments by those of $\eta$; polynomial factors are absorbed using
a slightly larger exponent $\theta<1$. Stirling's formula gives
\[
 \gamma_{m+1}(my)=\frac{e^{-mI(y)}}{\sqrt{2\pi m}}
                       \left(1-\frac1{12m}+O(m^{-2})\right).
\]
Multiplication proves \eqref{eq:saddle}. Positivity of $\HH$ on compacts
justifies the relative error and its logarithm.
\end{proof}

\begin{theorem}[Clipped profile expansion]
\label{thm:transfer}
In addition to \cref{lem:saddle}, suppose each $g_m$ is continuous and
log-concave on $(0,\infty)$, tends to zero at both endpoints, and has
positive interior values. For $z$ in a fixed compact subset of $(0,\infty)$,
put
\[
 a=(2\pi m)^{-1/2}e^{-mz},\qquad
 J_m(a)=\int_0^\infty\min\{1,g_m(u)/a\}\dd u.
\]
Then, uniformly in that set of $z$,
\begin{equation}
 J_m(a)=m(y_+-y_-)+[F]_{y_-}^{y_+}
                     +\frac1m[B]_{y_-}^{y_+}+O(m^{-2}).
 \label{eq:J}
\end{equation}
\end{theorem}
\begin{proof}
All roots remain in compact subsets of $(0,1)$ and $(1,\infty)$.
Consequently $s(y_\pm)$ are uniformly separated from zero.
Log-concavity makes the superlevel set $\{g_m\ge a\}$ an interval;
\cref{lem:saddle} places its two endpoints near $my_\pm$. A nonconstant
log-concave density cannot have a flat segment at a level below its maximum:
a flat segment would be a set of maximizers. Thus its endpoints are the
unique crossings, denoted $u_\pm$.

Taking logarithms in \eqref{eq:saddle} gives
\[
 \log(g_m(my)/a)=m(z-I(y))+L(y)+A(y)/m+O(m^{-2}).
\]
First this implies $u_\pm=my_\pm+O(1)$. Successive Taylor expansion then
gives
\begin{align}
 u_\pm&=my_\pm+b(y_\pm)+\frac{c(y_\pm)}m+O(m^{-2}),\label{eq:crossings}\\
 b(y)&=\frac{L(y)}{s(y)},\qquad
 c(y)=\frac{A(y)+L'(y)b(y)-b(y)^2/(2y^2)}{s(y)}.
 \label{eq:bc}
\end{align}
These conclusions need no derivative estimate on the saddle remainder:
the uniform $O(m^{-2})$ error and the nonzero slope bracket each crossing
within $O(m^{-2})$ of its candidate position.

We next retain both exterior tails. If $v$ stays in a compact subset of
$(1,\infty)$, integration by parts in the saddle formula gives
\begin{equation}
 \frac{\int_{mv}^{\infty}g_m(u)\dd u}{g_m(mv)}
 =\frac1{s(v)}+\frac1m
   \left(\frac{L'(v)}{s(v)^2}-\frac1{v^2s(v)^3}\right)+O(m^{-2}).
 \label{eq:righttail}
\end{equation}
On a compact subset of $(0,1)$ the corresponding lower-tail formula is
the negative of the right-hand expression:
\begin{equation}
 \frac{\int_0^{mv}g_m(u)\dd u}{g_m(mv)}
 =-\frac1{s(v)}-\frac1m
   \left(\frac{L'(v)}{s(v)^2}-\frac1{v^2s(v)^3}\right)+O(m^{-2}).
 \label{eq:lefttail}
\end{equation}
Here is a justification including the remote tails. The moment-generating
function is $(1-s)^{-(m+1)}\int e^{su}\eta_m(\dd u)$ for real $s<1$.
Markov's inequality, applied to the actual nonnegative law $g_m$, bounds
the upper tail beyond $mM$ by $C_Me^{-mI(M)}$, and the lower tail below
$m\delta$ by $C_\delta e^{-mI(\delta)}$.
Choose $M$ large and $\delta$ small so that these exponents exceed the
largest relevant endpoint exponent by a fixed positive amount. Both remote
tails, divided by the endpoint density, are exponentially negligible.
On the remaining compact intervals, \eqref{eq:saddle} has relative
remainder $O(m^{-2})$, whose integral is $O(m^{-2})$ times the leading
tail integral. For its smooth main part, integrate by parts twice using
$I'=s$, $I''=1/y^2$. The endpoint terms are
$1/s$ and $(L'/s^2-I''/s^3)/m$. A third integration, or an absolute bound
on the remaining derivative integral, gives the uniform $O(m^{-2})$
remainder. The amplitude $A/m$ cancels with its endpoint normalization
to this order. This proves both tail formulas, with the lower endpoint
orientation supplying the minus sign in \eqref{eq:lefttail}.

At $v=u_\pm/m$, substitute \eqref{eq:crossings} and use
$g_m(u_\pm)=a$. In a signed endpoint notation, the tail contribution is
\[
 \left[\frac1{s(y)}+\frac1m\left\{
 \frac{L'(y)-b(y)/y^2}{s(y)^2}-\frac1{y^2s(y)^3}
 \right\}\right]_{y_-}^{y_+}+O(m^{-2}).
\]
Adding the interval length $u_+-u_-$ gives \eqref{eq:J}, because
$b+1/s=F$ and
\[
 c+\frac{L'-b/y^2}{s^2}-\frac1{y^2s^3}
 =\frac A s+\frac{(L+1)L'}{s^2}
                 -\frac{L^2/2+L+1}{y^2s^3}=B.
\]
\end{proof}

\section{Verification for the geometric series}
\subsection{Exact signed convolution and its bounds}
For $a>0$ set
\begin{equation}
 \nu_a=\frac{\delta_0-e^{-a}\delta_a}{1-e^{-a}}.
 \label{eq:nu}
\end{equation}
The identity $p_a=\gamma_1*\nu_a$ follows by direct substitution.
The signed measure is an analytic device, never the law of a random
coordinate. For any subset of indices and any fixed $\theta<1$,
\begin{equation}
 \left\|\mathop{*}_{r\in D}\nu_{a_r}\right\|_\theta
 \le\prod_{r\ge1}\frac{1+e^{-(1-\theta)a_r}}{1-e^{-a_r}}<\infty.
 \label{eq:weighted}
\end{equation}
The same bound holds for $\theta<0$, using support in $[0,\infty)$.
Since $a_r$ grows geometrically, deleting the factors with $r>N$ changes
the product by $O(\sum_{r>N}e^{-(1-\theta)a_r})$ in this norm.
This error is smaller than every inverse power of $N$.

Let $C^j_{n,m}$ be the sum of the hidden coordinates under $Q_n$ and
let $E$ be an independent rate-one exponential. Its convolution density
$g^j_{n,m}$, the density of $C^j_{n,m}+E$, satisfies
\begin{align}
 g^{\mathrm{late}}_{n,m}
   &=\gamma_{m+1}*\eta_m^{\mathrm{late}},&
 \eta_m^{\mathrm{late}}
   &=\Law(R_0)*\mathop{*}_{r=1}^m\nu_{a_r},\label{eq:lateeta}\\
 g^{\mathrm{early}}_{n,m}
   &=\gamma_{m+1}*\eta_{n,m}^{\mathrm{early}},&
 \eta_{n,m}^{\mathrm{early}}
   &=\Law(R_0)*\mathop{*}_{r=n-m+1}^n\nu_{a_r}.
 \label{eq:earlyeta}
\end{align}
For the late mask the limiting signed transform is \eqref{eq:HL}.
For the early mask it is \eqref{eq:HE}; the convergence error is
superexponentially small in $n-m$. Both satisfy the weighted hypotheses
of \cref{lem:saddle}, uniformly in the critical window. The transfer
theorem permits such a two-index family whenever its bounds are uniform.

Each $p_a$ and the exponential density is log-concave. Finite convolution
preserves log-concavity \cite{prekopa}. Infinite convolution here follows
by weak limits; alternatively, convolve a finite partial sum with $E$,
pass to the continuous limiting densities on compact positive intervals,
and pass the defining log-concavity inequality to the limit.
The resulting density is positive for $u>0$, vanishes as $u\downarrow0$
and as $u\to\infty$, and is continuous. Positivity follows by making a
finite set of positive-cap coordinates small and bounding the remaining
deterministic tail; continuity follows by convolution with two bounded
continuous convolution factors. At the left endpoint,
$g(u)\le e^{-u}e^uP(C\le u)\to0$, since $C$ has no atom at zero.
Thus all hypotheses of \cref{thm:transfer} hold.

\subsection{The exact overlap identity}
Let $W^j_{n,m}$ be the observed-coordinate sum, let $f^j_{n,m}$ be the
density of $W^j_{n,m}-\E_{Q_n}W^j_{n,m}$, and let
$r^j_{n,m}=\E_{Q_n}C^j_{n,m}$. Define the full normalizer
\begin{equation}
 a_n=\E_{Q_n}\!\left[e^{T_n-\mu_n}\ind_{\{T_n\le\mu_n\}}\right].
\end{equation}
Conditioning the full likelihood from \eqref{eq:model} on the observed
vector gives its exact marginal likelihood
\begin{equation}
 \frac{g^j_{n,m}(r^j_{n,m}-v)}{a_n},
 \qquad v=W^j_{n,m}-\E_{Q_n}W^j_{n,m}.
\end{equation}
Indeed $e^{-u}\E[e^{C}\ind_{\{C\le u\}}]$ is the density of $C+E$.
It follows that
\begin{equation}
 \OO^j_{n,m}=\int_0^\infty f^j_{n,m}(r^j_{n,m}-u)
           \min\{1,g^j_{n,m}(u)/a_n\}\dd u.
 \label{eq:exactoverlap}
\end{equation}
This concerns the full vector marginal, not only an approximation of its sum.

\begin{lemma}[Uniform flattening in the critical window]
\label[lemma]{lem:flatten}
For $c_0\log n\le m\le c_1\log n$ and either mask,
\begin{equation}
 \frac{\OO^j_{n,m}}{\eps_n}
  =\int_0^\infty\min\{1,g^j_{n,m}(u)/\eps_n\}\dd u
                         +O(m^3/n).
 \label{eq:flatten}
\end{equation}
\end{lemma}
\begin{proof}
The observed sum has the exact density $\gamma_{n-m}*\xi^j_{n,m}$,
where $\xi^{\mathrm{late}}=\mathop{*}_{r=m+1}^n\nu_{a_r}$ and
$\xi^{\mathrm{early}}=\mathop{*}_{r=1}^{n-m}\nu_{a_r}$.
The weighted variations, all polynomial moments, and signed mean shifts
are uniformly bounded by \eqref{eq:weighted}. Stirling's formula and
Taylor expansion at the Gamma mean, integrated against these measures,
therefore give, for every fixed $M$,
\begin{equation}
 \sup_{|v|\le Mm}\left|\frac{f^j_{n,m}(v)}{\eps_n}-1\right|
       =O(m^2/n),\qquad \sup_v f^j_{n,m}(v)\le C\eps_n.
 \label{eq:local}
\end{equation}
Here the integral can first be restricted to signed shifts
$u\le M_1\log n$, with arbitrarily small inverse-polynomial error;
on the remaining interval the elementary bound is
$\sqrt{2\pi k}\gamma_k(k+w)=1+O((1+w^2)/k)$,
where $k=n-m$ and $w=O(\log n)$. This proves the stated error for
growing $m$, rather than invoking a fixed-$m$ expansion.

The full normalizer is independent of the observation mask and satisfies
\begin{equation}
 a_n=\eps_n(1+O(n^{-1})).\label{eq:normalizer}
\end{equation}
To see this directly, use the same signed representation with all $n$
bulk variables and the fixed boundary $R_0$. The centered total density,
on $|v|\le M_1\log n$, has the local expansion
$\eps_n[1+O((1+v^2)/n)+O((1+|v|^6)/n^2)]$ after integration of
the bounded signed moments. Integrating this density at $v=-e$ against
$e^{-e}\dd e$, $e\ge0$, proves \eqref{eq:normalizer}; its remote tail
is controlled using the global density bound. Also $r^j_{n,m}=m+O(1)$.

Choose constants $0<\delta<1<M$ so that $I(\delta)$ and $I(M)$ are
larger than any prescribed constant. The Chernoff estimates from the
proof of \cref{thm:transfer} imply
\[
 P(C^j_{n,m}+E\notin[\delta m,Mm])\le C e^{-Km}
\]
with $K$ arbitrarily large by this choice. Since $m\ge c_0\log n$,
the contributions outside this window, divided by $\eps_n$, are smaller
than $n^{-2}$ if $K$ is chosen sufficiently large. This uses the global
bound in \eqref{eq:local} and $a_n\asymp\eps_n$.
Inside the window, its length is $O(m)$ and \eqref{eq:local} gives a
total scaled flattening error $O(m^3/n)$.

Finally changing the clipping level from $a_n$ to $\eps_n$ costs
$O(m/n)$. In fact for $|\tau|\le1/2$,
\[
 |\min(1,x/(1+\tau))-\min(1,x)|
            \le2|\tau|\min(1,x),\qquad x\ge0,
\]
and the clipped integral is $O(m)$ by the same window and tail bounds.
This proves \eqref{eq:flatten}.
\end{proof}

\begin{proof}[Proof of \cref{thm:main}]
Apply \cref{thm:transfer} to \eqref{eq:lateeta} or \eqref{eq:earlyeta}.
Since $z=\log(n/m)/(2m)$ stays in a compact subset of $(0,\infty)$,
its uniformity conditions hold. Moreover
$(2\pi m)^{-1/2}e^{-mz}=\eps_n$ exactly.
The weighted convergence errors are smaller than $m^{-2}$ uniformly,
and \cref{lem:flatten} supplies the remaining $O(m^3/n)$ error.
\end{proof}

\section{Checks and limitations}
Two exact structural checks help detect errors in the coefficients.
For $\eta=\delta_b$, $b\ge0$, translation of the Gamma density leaves
$J_m$ unchanged. Substitution of $\HH(s)=e^{bs}$ into the formulas makes
the $b$ dependence cancel from both $[F]$ and $[B]$.
For $\eta=\delta_0$, $\HH=1$ and
\begin{equation}
 F=1/s,\qquad B=-\frac1{12s}-\frac1{y^2s^3}.
\end{equation}
The profile is then exactly Gamma$(m+1,1)$, so crossings and tails can
be evaluated independently with incomplete Gamma functions. The source
archive includes a high-precision diagnostic for this case and the
Gamma$(m+2,1)$ case $\HH(s)=(1-s)^{-1}$.
These floating-point evaluations are checks of the formulas, not interval
certificates or substitutes for the proof.

The theorem fixes $q$ and $\rho$. It does not give a uniform result as
$q\uparrow1$, as $c\downarrow0$, or as $c\uparrow\infty$. At the latter
boundary the endpoint slopes approach zero and the present endpoint
integration argument requires a different scaling. It does not solve
the entire early-mask range $m=o(n)$, arbitrary masks, general weights,
or the endpoint transseries for a fixed complement. It concerns overlap;
no new entropy theorem is asserted.

All new results here remain unformalized. No exact Lean declaration is
claimed as a counterpart to \cref{thm:main,thm:transfer,cor:comparison}.
The required formal obligations include measure-theoretic likelihood
reduction, signed convolution estimates, uniform asymptotic remainders,
and level-crossing analysis. The two source articles are manuscript
sources, not externally refereed results. Their exact algebra has been
rederived here in the scope needed for the new statements.

\section{Concrete next questions}
Three bounded extensions now look natural.
First, identify how the correction matches the thin-complement endpoint
when $m/\log n\to0$, where $y_-$ approaches zero and the signed moment
bounds must be used on a growing range of tilt parameters.
Second, obtain a transition theorem as $m/\log n\to\infty$ by replacing
the separated-endpoint expansion with a uniform central saddle analysis.
Third, classify arbitrary masks using their limiting correction transforms.
The strict ordering proved above suggests studying whether a general
ordering of the hidden-cap profiles induces an ordering of the normalized
secant correction. Each extension requires new estimates; none is implicit
in the present compact-critical theorem.

\begin{thebibliography}{9}
\bibitem{repo-late}
ProveIt research manuscript collection, maintained by V.~Reshetnikov.
\emph{Critical Complements in Fabius Conditioning: Sharp Information Loss,
Geometric Phase Constants, and a Logarithmic Crossover}.
Snapshot inspected 1 October 2026, commit
\nolinkurl{63a7a325109ba611a1816b61dfd0eb072b896a7a}.
\href{https://github.com/VladimirReshetnikov/ProveIt/blob/63a7a325109ba611a1816b61dfd0eb072b896a7a/Analysis/FabiusFunction/docs/semi-formalized-research-frontiers/drafts/representations/Critical_Complements_Sharp_Information_Loss/article.tex}{Pinned manuscript source}.
Source of the leading crossover and the explicitly open second-order question.

\bibitem{repo-early}
ProveIt research manuscript collection, maintained by V.~Reshetnikov.
\emph{Critical Complements in Fabius Conditioning: Logarithmic Overlap,
Finite-Memory Corrections, and the Order-Zero R\'enyi Crossover}.
Same inspected snapshot.
\href{https://github.com/VladimirReshetnikov/ProveIt/blob/63a7a325109ba611a1816b61dfd0eb072b896a7a/Analysis/FabiusFunction/docs/semi-formalized-research-frontiers/drafts/representations/Critical_Complements_Fabius_Conditioning/article.tex}{Pinned manuscript source}.
Source of the distinct early-hidden mask and its growing-complement question.

\bibitem{arias}
J.~Arias de Reyna.
\emph{An Infinitely Differentiable Function with Compact Support:
Definition and Properties}.
English translation of the 1982 paper, arXiv:1702.05442 (2017).
\url{https://arxiv.org/abs/1702.05442}.
Background for the probability interpretation of the Fabius--Rvachev function.

\bibitem{fmnn}
V.~F\'eray, P.-L.~M\'eliot and A.~Nikeghbali.
\emph{Mod-$\phi$ Convergence I: Normality Zones and Precise Deviations}.
arXiv:1304.2934v4 (2015).
\url{https://arxiv.org/abs/1304.2934}.
General precise-deviation context; no result from this source is used as an
unverified uniform estimate in the proof.

\bibitem{dlmf}
NIST Digital Library of Mathematical Functions, Section 8.11(iii).
\emph{Incomplete Gamma Functions: Large Parameter with Fixed Ratio}.
\url{https://dlmf.nist.gov/8.11.iii}.
Classical Gamma-tail asymptotic context.

\bibitem{prekopa}
A.~Pr\'ekopa.
\emph{On Logarithmic Concave Measures and Functions}.
Acta Scientiarum Mathematicarum (Szeged) \textbf{34} (1973), 335--343.
\url{https://rutcor.rutgers.edu/~prekopa/SCIENT2.pdf}.
Convolution preservation of log-concavity, Theorem 7.
\end{thebibliography}
\end{document}
